% GENERATED FILE. Edit canonical lessons, then run npm run generate:books.
% canonical-source-hash: fnv1a64:8f3591a19c220274
% canonical-lessons: SA-C123-pitham, SA-C123-dvicakrika, SA-C123-yanam, SA-C123-nauyanam, SA-C123-dhumayanam

\chapter{Seat, Bicycle, Vehicle, Ship, Train}
\label{ch:sa-seat-bicycle-vehicle-ship-train}
\hlchaptermodality{123}

\begin{chapteropening}
\textbf{By the end of this chapter:} \emph{I can say a seat, a bicycle, a vehicle, a ship and a train.}

Complete the last lesson of chapter 123: I can say a seat, a bicycle, a vehicle, a ship and a train.
\end{chapteropening}

\section[\texorpdfstring{\sk{पीठम्} (pīṭham)}{पीठम् (pīṭham)}]{\sk{पीठम्} (pīṭham) — a seat, a stool}

\label{lesson:SA-C123-pitham}



\begin{quote}
Before the new one: say the Sanskrit for a pot, then the Sanskrit for a cot.
\end{quote}

\subsection*{You'll want to know: \sk{पीठम्}}
\textbf{\sk{पीठम्}} — \emph{pīṭham} — \textquotedblleft{}a seat, a stool\textquotedblright{}.

\begin{cousinweb}
\emph{pīṭha} \textquotedblleft{}a stool, a seat\textquotedblright{} has an uncertain origin: MW wonders whether it is a worn form of \emph{pi-sad} \textquotedblleft{}sit upon\textquotedblright{}, while CDIAL 8222 thinks the variety of its forms points to a non-Aryan source. It lives on as Hindi \emph{pīṛhā} \textquotedblleft{}low stool\textquotedblright{} (CDIAL 8222).
\end{cousinweb}

\begin{grammarlens}[title={What you've built}]
One more word for this chapter.
\end{grammarlens}

\subsection*{Your turn}
\begin{itemize}
\raggedright
  \item \emph{Say it:} \emph{pīṭham}
  \item \emph{Say it:} \emph{pīṭham}, once more
  \item \emph{Say it:} say \emph{pīṭham}
  \item \emph{Recall:} say the Sanskrit for a pot, then the Sanskrit for a cot, then say \emph{pīṭham} again
\end{itemize}

\begin{tcolorbox}[breakable,colback=teal!4,colframe=teal!35!black,title={Before you move on}]
What does \sk{पीठम्} mean? (\textquotedblleft{}A seat, a stool\textquotedblright{}.) Say it once more.
\end{tcolorbox}
\section[\texorpdfstring{\sk{द्विचक्रिका} (dvicakrikā)}{द्विचक्रिका (dvicakrikā)}]{\sk{द्विचक्रिका} (dvicakrikā) — a bicycle}

\label{lesson:SA-C123-dvicakrika}



\begin{quote}
Before the new one: say the Sanskrit for a cot, then the Sanskrit for a seat.
\end{quote}

\subsection*{You'll want to know: \sk{द्विचक्रिका}}
\textbf{\sk{द्विचक्रिका}} — \emph{dvicakrikā} — \textquotedblleft{}a bicycle\textquotedblright{}.

\begin{cousinweb}
\emph{dvicakrikā} is built from \emph{dvi} \textquotedblleft{}two\textquotedblright{} and \emph{cakra} \textquotedblleft{}wheel\textquotedblright{} with the suffix \emph{-ikā}: \textquotedblleft{}the little two-wheeler\textquotedblright{}. It is a modern coinage, made the way English \emph{bicycle} was made from Latin \emph{bi-} \textquotedblleft{}two\textquotedblright{} and Greek \emph{kyklos} \textquotedblleft{}wheel\textquotedblright{}; \emph{cakra} lives on as Marathi \emph{cāk} \textquotedblleft{}wheel\textquotedblright{} (CDIAL 4538).
\end{cousinweb}

\begin{grammarlens}[title={What you've built}]
One more word for this chapter.
\end{grammarlens}

\subsection*{Your turn}
\begin{itemize}
\raggedright
  \item \emph{Say it:} \emph{dvicakrikā}
  \item \emph{Say it:} \emph{dvicakrikā}, once more
  \item \emph{Say it:} say \emph{dvicakrikā}
  \item \emph{Recall:} say the Sanskrit for a cot, then the Sanskrit for a seat, then say \emph{dvicakrikā} again
\end{itemize}

\begin{tcolorbox}[breakable,colback=teal!4,colframe=teal!35!black,title={Before you move on}]
What does \sk{द्विचक्रिका} mean? (\textquotedblleft{}A bicycle\textquotedblright{}.) Say it once more.
\end{tcolorbox}
\section[\texorpdfstring{\sk{यानम्} (yānam)}{यानम् (yānam)}]{\sk{यानम्} (yānam) — a vehicle}

\label{lesson:SA-C123-yanam}



\begin{quote}
Before the new one: say the Sanskrit for a seat, then the Sanskrit for a bicycle.
\end{quote}

\subsection*{You'll want to know: \sk{यानम्}}
\textbf{\sk{यानम्}} — \emph{yānam} — \textquotedblleft{}a vehicle\textquotedblright{}.

\begin{cousinweb}
\emph{yāna} comes from the root √yā \textquotedblleft{}go\textquotedblright{} (CDIAL 10459): first \textquotedblleft{}a journey, going\textquotedblright{}, then \textquotedblleft{}a vehicle of any kind\textquotedblright{}, as in \emph{mahāyāna} \textquotedblleft{}the great vehicle\textquotedblright{} (MW). The root is inherited from Indo-European (EtymDB).
\end{cousinweb}

\begin{grammarlens}[title={What you've built}]
One more word for this chapter.
\end{grammarlens}

\subsection*{Your turn}
\begin{itemize}
\raggedright
  \item \emph{Say it:} \emph{yānam}
  \item \emph{Say it:} \emph{yānam}, once more
  \item \emph{Say it:} say \emph{yānam}
  \item \emph{Recall:} say the Sanskrit for a seat, then the Sanskrit for a bicycle, then say \emph{yānam} again
\end{itemize}

\begin{tcolorbox}[breakable,colback=teal!4,colframe=teal!35!black,title={Before you move on}]
What does \sk{यानम्} mean? (\textquotedblleft{}A vehicle\textquotedblright{}.) Say it once more.
\end{tcolorbox}
\section[\texorpdfstring{\sk{नौयानम्} (nauyānam)}{नौयानम् (nauyānam)}]{\sk{नौयानम्} (nauyānam) — a ship}

\label{lesson:SA-C123-nauyanam}



\begin{quote}
Before the new one: say the Sanskrit for a bicycle, then the Sanskrit for a vehicle.
\end{quote}

\subsection*{You'll want to know: \sk{नौयानम्}}
\textbf{\sk{नौयानम्}} — \emph{nauyānam} — \textquotedblleft{}a ship\textquotedblright{}.

\begin{cousinweb}
\emph{nauyāna} is built from \emph{nau} \textquotedblleft{}ship, boat\textquotedblright{} and \emph{yāna} \textquotedblleft{}going, vehicle\textquotedblright{}, from the root √yā \textquotedblleft{}go\textquotedblright{}; MW already records it as \textquotedblleft{}a ship\textquotedblright{}. MW sets \emph{nau} beside Greek \emph{naus} and Latin \emph{navis}, the source of English \emph{navy} and \emph{naval}.
\end{cousinweb}

\begin{grammarlens}[title={What you've built}]
One more word for this chapter.
\end{grammarlens}

\subsection*{Your turn}
\begin{itemize}
\raggedright
  \item \emph{Say it:} \emph{nauyānam}
  \item \emph{Say it:} \emph{nauyānam}, once more
  \item \emph{Say it:} say \emph{nauyānam}
  \item \emph{Recall:} say the Sanskrit for a bicycle, then the Sanskrit for a vehicle, then say \emph{nauyānam} again
\end{itemize}

\begin{tcolorbox}[breakable,colback=teal!4,colframe=teal!35!black,title={Before you move on}]
What does \sk{नौयानम्} mean? (\textquotedblleft{}A ship\textquotedblright{}.) Say it once more.
\end{tcolorbox}
\section[\texorpdfstring{\sk{धूमयानम्} (dhūmayānam)}{धूमयानम् (dhūmayānam)}]{\sk{धूमयानम्} (dhūmayānam) — a train}

\label{lesson:SA-C123-dhumayanam}



\begin{quote}
Before the new one: say the Sanskrit for a vehicle, then the Sanskrit for a ship.
\end{quote}

\subsection*{You'll want to know: \sk{धूमयानम्}}
\textbf{\sk{धूमयानम्}} — \emph{dhūmayānam} — \textquotedblleft{}a train\textquotedblright{}.

\begin{cousinweb}
\emph{dhūmayāna} is built from \emph{dhūma} \textquotedblleft{}smoke\textquotedblright{} and \emph{yāna} \textquotedblleft{}vehicle\textquotedblright{}: \textquotedblleft{}the smoke-carriage\textquotedblright{}, a modern coinage from the days of the steam engine. MW sets \emph{dhūma} beside Latin \emph{fumus} \textquotedblleft{}smoke\textquotedblright{}, the source of English \emph{fume}, and it lives on as Hindi \emph{dhuāṁ} \textquotedblleft{}smoke\textquotedblright{} (CDIAL 6849).
\end{cousinweb}

\begin{grammarlens}[title={What you've built}]
One more word for this chapter.
\end{grammarlens}

\subsection*{Your turn}
\begin{itemize}
\raggedright
  \item \emph{Say it:} \emph{dhūmayānam}
  \item \emph{Say it:} \emph{dhūmayānam}, once more
  \item \emph{Say it:} say \emph{dhūmayānam}
  \item \emph{Recall:} say the Sanskrit for a vehicle, then the Sanskrit for a ship, then say \emph{dhūmayānam} again
\end{itemize}

\begin{tcolorbox}[breakable,colback=teal!4,colframe=teal!35!black,title={Before you move on}]
What does \sk{धूमयानम्} mean? (\textquotedblleft{}A train\textquotedblright{}.) Say it once more.
\end{tcolorbox}
